\chaptertitle{第四章\quad 排列}

第三章我们掌握了"分步相乘"这个基本功。现在来研究一个非常经典的问题：从$n$个不同的元素中选出$m$个，并按一定的顺序排成一列，有多少种排法？这种排法就叫\textbf{排列}。

\sectiontitle{4.1\quad 全排列}

假设$3$个人排成一排照相，有多少种不同的站法？

把$3$个位置排好：第一个位置$3$人中任选一个（$3$种），第二个位置从剩下$2$人中选（$2$种），第三个位置只剩$1$人（$1$种）：
\[
3\times2\times1 = 6\ \text{种}
\]

\begin{center}
\begin{tikzpicture}[scale=0.5]
\node at (0,3) {甲、乙、丙排成一排的$6$种可能：};
\node at (0,1) {甲\;乙\;丙};
\node at (0,0) {甲\;丙\;乙};
\node at (0,-1) {乙\;甲\;丙};
\node at (0,-2) {乙\;丙\;甲};
\node at (0,-3) {丙\;甲\;乙};
\node at (0,-4) {丙\;乙\;甲};
\end{tikzpicture}
\end{center}

\textbf{全排列}：将$n$个不同元素按一定顺序全部排成一列，共有
\[
n! = n\times(n-1)\times(n-2)\times\cdots\times2\times1
\]
种排法，读作"$n$的阶乘"。

\begin{example}
计算：(1) $4!$ \qquad (2) $5!$ \qquad (3) $1!$
\end{example}
\begin{solution}
(1) $4! = 4\times3\times2\times1 = 24$\\
(2) $5! = 5\times4\times3\times2\times1 = 120$\\
(3) $1! = 1$
\end{solution}

\begin{exercise}
\item 计算$6!$。
\item $5$个人排成一排照相，有多少种站法？
\item $4$本不同的书放在书架上排成一排，有多少种排法？
\end{exercise}

\sectiontitle{4.2\quad 选排列}

如果不需要全部元素都排，只从$n$个不同元素中选出$m$个来排成一列（$m\le n$），这就是\textbf{选排列}，记作$A_n^m$（有些书也记作$P_n^m$）。

\begin{example}
从$5$个人中选$3$人排成一排（$3$个位置），有多少种排法？
\end{example}
\begin{solution}
第一个位置$5$种，第二个位置$4$种，第三个位置$3$种：
$A_5^3 = 5\times4\times3 = 60$种。
\end{solution}

几个具体例子，观察模式：
\[
A_5^2 = 5\times4 = 20,\qquad
A_5^3 = 5\times4\times3 = 60,\qquad
A_5^4 = 5\times4\times3\times2 = 120
\]

\begin{example}
从这些例子中，你能归纳出$A_n^m$的一般公式吗？
\end{example}
\begin{solution}
从$n$个中选$m$个排列，相当于从$n$开始往下乘$m$个数：
\[
A_n^m = n\times(n-1)\times(n-2)\times\cdots\times(n-m+1) = \frac{n!}{(n-m)!}
\]
最后一步是因为$n! = n\times(n-1)\times\cdots\times(n-m+1)\times(n-m)!$，所以
\[
A_n^m = \frac{n!}{(n-m)!}
\]
当$m=n$时，$A_n^n = n!$，与全排列一致。
\end{solution}

\begin{example}
从$8$人中选$3$人站成一排领奖（分一、二、三等奖），有多少种可能？
\end{example}
\begin{solution}
$A_8^3 = 8\times7\times6 = 336$种。
\end{solution}

\begin{exercise}
\item 计算：$A_7^2$，$A_6^3$，$A_{10}^4$。
\item 从$10$本书中选$3$本按顺序放在书架上，有多少种放法？
\item $6$个人中选$4$人排成一排，有多少种排法？
\end{exercise}

\sectiontitle{4.3\quad 有限重复元素的排列}

如果元素不是全都不同呢？比如把"APPLE"这个单词的字母排成一排，有多少种不同的排法？

字母$A,P,P,L,E$中$P$出现了两次。如果它们全是不同的，答案是$5! = 120$。但两个$P$是一样的，交换它们不产生新排列。

\begin{example}
计算"APPLE"的字母能组成多少个不同的排列。
\end{example}
\begin{solution}
把$5$个位置分配字母：
先写$5!$种（假设所有字母不同），再除以两个$P$互换带来的重复倍数$2!$：
\[
\frac{5!}{2!} = \frac{120}{2} = 60
\]
\end{solution}

更一般地：如果$n$个元素中，有$n_1$个相同、$n_2$个相同、……、$n_k$个相同（$n_1+n_2+\cdots+n_k=n$），则排列数为
\[
\frac{n!}{n_1!\,n_2!\,\cdots\,n_k!}
\]

\begin{example}
"MATHEMATICS"这个单词的字母能组成多少个不同的排列？
\end{example}
\begin{solution}
字母：$M\times2$、$A\times2$、$T\times2$、$H$、$E$、$I$、$C$、$S$，共$11$个。
\[
\frac{11!}{2!\,2!\,2!} = \frac{39916800}{8} = 4989600
\]
\end{solution}

\begin{exercise}
\item "BOOK"的字母能组成多少个不同的排列？
\item "MISSISSIPPI"的字母能组成多少个不同的排列？
\end{exercise}

\sectiontitle{4.4\quad 圆排列}

前面的排列都是站成一排（线排列）。如果围成一个圆桌坐呢？

$4$个人围圆桌而坐，有多少种不同的坐法？

关键区别：在线排列中，$ABCD$和$BCDA$是不同的（因为位置不同）。但在圆排列中，如果整体旋转后重合，算同一种坐法。

\begin{example}
$4$个人围圆桌而坐，有多少种不同的坐法？
\end{example}
\begin{solution}
先固定其中一人（比如甲）的位置——因为圆桌旋转等价，固定一个人就消除了旋转重复。
剩下的$3$个人在剩下的$3$个位置上全排列：$3! = 6$种。

一般的\textbf{圆排列}公式：$n$个不同元素围成一圈，有
\[
\frac{n!}{n} = (n-1)!
\]
种排列方式。
\end{solution}

\begin{example}
$6$颗不同的珍珠穿成一个手链（可以翻转），有多少种穿法？
\end{example}
\begin{solution}
先考虑圆排列：$(6-1)! = 120$种。\\
但手链可以翻转——翻过来之后$ABCDEF$变成$AFEDCB$。在圆排列中它们本来算两种，但在手链中它们是一样的（从不同面看）。\\
所以再除以$2$：$120\div2 = 60$种。
\end{solution}

\begin{exercise}
\item $5$个人围圆桌开会，有多少种座位安排？
\item $8$个人围圆桌而坐，其中两人要求坐在一起，有多少种坐法？
\end{exercise}

\sectiontitle{4.5\quad 可重复排列}

如果元素可以重复选呢？比如一个$3$位数的密码锁，每位从$0$到$9$中选，可以重复。

这就是\textbf{可重复排列}：从$n$个不同元素中取$m$个排成一列（每个元素可以重复取），共有
\[
n^m
\]
种排列。

\begin{example}
(1) 一个$3$位数的密码锁有多少种可能？\\
(2) 一个$4$位数的手机密码（每位$0$-$9$）有多少种可能？\\
(3) 抛$5$次硬币（正/反），有多少种结果序列？
\end{example}
\begin{solution}
(1) $10^3 = 1000$种。\\
(2) $10^4 = 10000$种。\\
(3) $2^5 = 32$种。
\end{solution}

\begin{exercise}
\item 一个$6$位数的验证码（每位$0$-$9$）有多少种可能？
\item 每个同学在$A,B,C,D$四个选项中选择一个答案，$10$个同学的答案序列有多少种？
\end{exercise}

\sectiontitle{4.6\quad 限制条件问题}

实际中的排列问题往往不是直接套公式——需要加上各种限制条件。

\begin{example}
$5$人排成一排，甲不站两端，有多少种排法？
\end{example}
\begin{solution}
方法一（从位置出发）：中间$3$个位置甲都可以站，所以甲有$3$种选择。剩下$4$人全排列$4! = 24$种。总数$3\times24=72$种。

方法二（用总数减去甲站两端的情况）：\\
$5$人全排列$5! = 120$种。\\
甲站左端：剩下$4$人全排列$24$种。\\
甲站右端：同理$24$种。\\
总数$120 - 24 - 24 = 72$种。
\end{solution}

\begin{example}
$4$名男生和$3$名女生排成一排，要求女生互不相邻，有多少种排法？
\end{example}
\begin{solution}
先排$4$名男生（不限制）：$4! = 24$种。\\
男生排好后产生$5$个空位（两端也算）：\;\_\_N\_\_N\_\_N\_\_N\_\_\\
从$5$个空位中选$3$个放女生（顺序重要）：$A_5^3 = 5\times4\times3 = 60$种。\\
总数$24\times60 = 1440$种。
\end{solution}

\begin{example}
$3$男$3$女排成一排，要求男女交替，有多少种排法？
\end{example}
\begin{solution}
两种可能：男男女女或女女男男。\\
情况$1$（男女男女男女）：男生在奇数位$3!$种，女生在偶数位$3!$种：$3!\times3! = 36$种。\\
情况$2$（女男女男女男）：同样$3!\times3! = 36$种。\\
总数$36+36=72$种。
\end{solution}

\begin{exercise}
\item $6$人排成一排，甲不站排头，乙不站排尾，有多少种排法？
\item $3$名男生和$2$名女生排成一排，要求女生不能相邻，有多少种排法？
\item $4$人排成一排，甲和乙必须相邻，有多少种排法？
\end{exercise}

\sectiontitle{本章小结}

本章学习了排列的各种形式：全排列（$n!$）、选排列（$A_n^m$）、有限重复元素的排列（$\dfrac{n!}{n_1!\,n_2!\,\cdots}$）、圆排列（$(n-1)!$）、可重复排列（$n^m$）。解决限制条件的排列问题，核心思路是：先满足限制条件，再排剩余元素；或者用总数减去不符合条件的情况（补集思想）。

\sectiontitle{课后练习}

\subsection*{A组\quad 基础巩固}

\begin{exercise}
\item 计算：$A_8^3$，$\dfrac{7!}{3!}$，$(5-1)!$。
\item $7$本书排成一排，有多少种排法？
\item 从$10$人中选$4$人站成一排领奖，有多少种可能？
\item "BANANA"的字母能组成多少个不同的排列？
\item $6$人排成一排，甲必须站在第一位，有多少种排法？
\item $8$人围圆桌开会，有多少种座位安排？
\end{exercise}

\subsection*{B组\quad 挑战提升}

\begin{exercise}
\item 用$0,1,2,3,4,5$能组成多少个无重复数字的六位数？其中有多少个是偶数？
\item $4$对夫妇排成一排，要求每对夫妇都不相邻，有多少种排法？（提示：先求总数，减去至少一对相邻的情况。）
\item "MATHEMATICS"中有多少个字母排列是元音字母（$A,E,I$）不相邻的？
\item $n$个人围圆桌而坐，其中甲和乙相对而坐，有多少种坐法？
\end{exercise}

\vfill
\noindent\rule{\textwidth}{0.4pt}
本章练习答案请参见书末附录。
